% ===== §7 =====
\section{The Quality of a Shared Life}
\label{sec:wellbeing}

\noindent
A fourth work of relational understanding concerns the quality of a shared life, what makes a bond good to be within and not merely stable. The established accounts here are several, and one may be taken as representative. Self-determination theory, developed by Deci and Ryan~\citep{deci2000self}, holds that wellbeing rests on the satisfaction of three basic psychological needs, for autonomy, for competence, and for relatedness, and it treats the quality of a relationship as bound up with the degree to which the bond satisfies the need for relatedness, the sense of being connected to and cared for by others. A large body of work on relationship satisfaction, taking various forms, likewise renders quality as a level of satisfaction, a favorable balance of rewards over costs, or a fulfillment of needs and expectations that the partners bring to the bond.

These accounts reach something and render it in a particular form. They locate the quality of a bond in the satisfaction of needs the individuals bring, and they treat that quality as a level, higher or lower, more or less sustained. Relatedness, on the self-determination account, is a need of the individual that the relationship meets, and quality is the degree of its meeting. This is a rendering of real importance, and it locates quality in the individual and the individual's needs, treating the relationship as the means by which a need belonging to each is satisfied.

The relational character supplies a different location for the quality of a bond. On the present account the quality of an intimate relationship is not, at its root, the degree to which it satisfies needs that the individuals bring to it. It is the richness of what is generated between them, the depth and the reach of the shared reality that their crossing of worlds produces. A bond is good to be within, on this account, to the degree that it generates, in the relation between two worlds, observations and values and a shared life that neither world could hold alone. This relocates quality from the satisfaction of antecedent individual needs to the generativity of the relation itself, and it accounts for something the need-satisfaction picture renders with difficulty, that the deepest bonds are often those in which the partners are most changed by the relation, most drawn beyond the needs and the selves they brought to it, which is hard to state as the efficient satisfaction of needs given in advance.

The generative character distinguishes a bond that is good in this sense from one that is merely stable and satisfied. A relationship can reach a stable level of need-satisfaction and rest there, and on the need account this is quality achieved. On the present account the quality of a bond is a generativity sustained and not a level reached and maintained, the continued production, in the crossing of the two worlds, of a shared reality that goes on growing. The distinction matters for the understanding of long bonds. A bond may be stable and satisfied and yet cease to generate, the partners settled into a fixed mutual rendering that no longer crosses into new observation, and on the present account such a bond has lost the quality that the need account, registering only its stability and satisfaction, would still record. The good bond is the generative bond, the one whose crossing of worlds continues to produce what neither world held before, and this is a matter of movement and not of a level, of a spiral that returns to the shared life and finds it grown where a circle would return it unchanged.

The relation to the received accounts is again one of location. The satisfaction of needs for autonomy, competence, and relatedness is real, and its measurement tracks something that matters. The present account holds that these render the individual conditions and the felt registers of a bond's quality, and that they leave for a relational and generative account the matter of what that quality consists in, which is the richness and the continued growth of what two worlds generate between them.
